\sectionheading{5.30 Python Program-এর প্রাথমিক পরিচিতি}

কোনো নির্দিষ্ট কাজ সম্পন্ন বা Problem সমাধানের জন্য Python
ভাষায় এক বা একাধিক Instruction বা Statement লিখে
Python Program তৈরি করা হয়।

%================================================
% Definition
%================================================

\booksubheading{Python Program কী?}

\begin{definitionbox}

\textbf{\color{darkTeal}সংজ্ঞা}

কোনো Problem সমাধান বা নির্দিষ্ট কাজ সম্পন্ন করার জন্য
Python Programming Language-এর নিয়ম অনুসারে লেখা
এক বা একাধিক \textbf{Statement-এর সমষ্টিকে Python Program}
বলে।

\end{definitionbox}

\textbf{সহজভাবে মনে রাখবে—}

\begin{center}

\textbf{Instruction / Statement}
\(\longrightarrow\)
\textbf{Program}
\(\longrightarrow\)
\textbf{Execution}
\(\longrightarrow\)
\textbf{Output}

\end{center}

%================================================
% Python Statement
%================================================

\booksubheading{Python Statement}

Python Program-এর প্রতিটি কার্যকর নির্দেশনাকে সাধারণভাবে
\textbf{Statement} বলা হয়।

\begin{examplebox}{উদাহরণ}

\begin{codebox}

\begin{english}
\ttfamily
print("Hello")
\end{english}

\end{codebox}

এখানে পুরো লাইনটি একটি Python Statement।

\medskip

আরও একটি উদাহরণ—

\begin{codebox}

\begin{english}
\ttfamily
print(10 + 20)
\end{english}

\end{codebox}

\textbf{Output:}

\begin{codebox}

\begin{english}
\ttfamily
30
\end{english}

\end{codebox}

\end{examplebox}

%================================================
% Simple Python Program
%================================================

\replaceableimage[68mm]
{images/python-editor-run.png}
{Python Editor-এ Program লিখে Run করার পর Output}

\begin{center}
\footnotesize\textbf{চিত্র: Python Editor-এ Program লিখে Run করার ধাপ}
\end{center}

\booksubheading{একটি সাধারণ Python Program}

\begin{codebox}

\begin{english}
\ttfamily
print("Hello, Python!")\\
print("Welcome to Programming")
\end{english}

\end{codebox}

\textbf{Output:}

\begin{codebox}

\begin{english}
\ttfamily
Hello, Python!\\
Welcome to Programming
\end{english}

\end{codebox}

এখানে দুটি Statement রয়েছে।

Python সাধারণভাবে Statementগুলো
\textbf{উপর থেকে নিচের দিকে (Top to Bottom)}
Execute করে।

\begin{center}

\textbf{১ম Statement}
\(\longrightarrow\)
\textbf{২য় Statement}
\(\longrightarrow\)
\textbf{৩য় Statement}
\(\longrightarrow\)
\textbf{\dots}

\end{center}

%================================================
% How Program Works
%================================================

\booksubheading{Python Program কীভাবে কাজ করে?}

একটি Python Program-এর সাধারণ কার্যধারা—

\begin{center}

\textbf{Source Code}
\(\longrightarrow\)
\textbf{Python Interpreter}
\(\longrightarrow\)
\textbf{Execution}
\(\longrightarrow\)
\textbf{Output}

\end{center}

Programmer প্রথমে Python Code লেখে।

Python Interpreter সেই Code Execute করে এবং প্রয়োজন অনুযায়ী
Result বা Output প্রদান করে।

%================================================
% Source Code
%================================================

\booksubheading{Source Code কী?}

\begin{definitionbox}

Programmer কোনো Programming Language ব্যবহার করে যে Code
লেখে তাকে \textbf{Source Code} বলা হয়।

\end{definitionbox}

উদাহরণ—

\begin{codebox}

\begin{english}
\ttfamily
print("ICT")
\end{english}

\end{codebox}

এটি Python-এর Source Code।

\begin{rememberbox}

\textbf{\color{deepOrange}মনে রাখবে}

\textbf{Programmer লিখে}
\(\longrightarrow\)
Source Code

\medskip

\textbf{Interpreter Execute করে}
\(\longrightarrow\)
Program চলে

\end{rememberbox}

%================================================
% Python Program File
%================================================

\booksubheading{Python Program File}

Script Mode-এ Python Program সাধারণত
\textbf{\textenglish{\ttfamily .py}}
Extension-যুক্ত File-এ Save করা হয়।

উদাহরণ—

\begin{codebox}

\begin{english}
\ttfamily
hello.py\\
first\_program.py
\end{english}

\end{codebox}

এরপর File Run করলে Python Program Execute হয়।

\begin{rememberbox}

\textbf{\color{deepOrange}MCQ তথ্য}

Python Source Program File Extension
\(\longrightarrow\)
\textbf{\textenglish{\ttfamily .py}}

\end{rememberbox}

%================================================
% print()
%================================================

\booksubheading{\textenglish{\ttfamily print()} সম্পর্কে প্রাথমিক ধারণা}

Python-এ Screen-এ Text, Number বা কোনো Result প্রদর্শনের জন্য
সাধারণত
\textenglish{\ttfamily print()}
Function ব্যবহার করা হয়।

\begin{examplebox}{Text প্রদর্শন}

\begin{codebox}

\begin{english}
\ttfamily
print("Bangladesh")
\end{english}

\end{codebox}

\textbf{Output:}

\begin{codebox}

\begin{english}
\ttfamily
Bangladesh
\end{english}

\end{codebox}

\end{examplebox}

\begin{examplebox}{Calculation-এর Result প্রদর্শন}

\begin{codebox}

\begin{english}
\ttfamily
print(5 + 3)
\end{english}

\end{codebox}

\textbf{Output:}

\begin{codebox}

\begin{english}
\ttfamily
8
\end{english}

\end{codebox}

\end{examplebox}

\begin{rememberbox}

\textbf{\color{deepOrange}আপাতত মনে রাখবে}

\begin{center}

\textenglish{\ttfamily print()}
\(\longrightarrow\)
\textbf{Output প্রদর্শন}

\end{center}

এর বিস্তারিত ব্যবহার
\textbf{Input ও Output}
অংশে আলোচনা করা হবে।

\end{rememberbox}

%================================================
% Text and Number
%================================================

\booksubheading{Text ও Number-এর প্রাথমিক ধারণা}

Python-এ Text সাধারণত Quotation Mark-এর মধ্যে লেখা হয়।

\begin{codebox}

\begin{english}
\ttfamily
print("Hello")
\end{english}

\end{codebox}

অন্যদিকে Number সরাসরি লেখা যায়—

\begin{codebox}

\begin{english}
\ttfamily
print(100)
\end{english}

\end{codebox}

\textbf{Output:}

\begin{codebox}

\begin{english}
\ttfamily
100
\end{english}

\end{codebox}

%================================================
% Quotation Difference
%================================================

\booksubheading{Quotation Mark থাকলে কী পার্থক্য হয়?}

নিচের দুটি Statement লক্ষ্য করো—

\begin{examplebox}{Quotation-এর মধ্যে}

\begin{codebox}

\begin{english}
\ttfamily
print("10 + 20")
\end{english}

\end{codebox}

\textbf{Output:}

\begin{codebox}

\begin{english}
\ttfamily
10 + 20
\end{english}

\end{codebox}

এখানে
\textenglish{\ttfamily 10 + 20}
Text হিসেবে প্রদর্শিত হয়েছে।

\end{examplebox}

\begin{examplebox}{Quotation ছাড়া}

\begin{codebox}

\begin{english}
\ttfamily
print(10 + 20)
\end{english}

\end{codebox}

\textbf{Output:}

\begin{codebox}

\begin{english}
\ttfamily
30
\end{english}

\end{codebox}

এখানে
\textenglish{\ttfamily 10 + 20}
একটি Calculation হিসেবে কাজ করেছে।

\end{examplebox}

\begin{rememberbox}

\textbf{\color{deepOrange}গুরুত্বপূর্ণ}

\textbf{Quotation-এর ভিতরে}
\(\longrightarrow\)
Text

\medskip

\textbf{Quotation ছাড়া}
\(\longrightarrow\)
Expression / Calculation

\end{rememberbox}

%================================================
% Top to Bottom
%================================================

\booksubheading{Python Program সাধারণভাবে Top to Bottom Execute হয়}

নিচের Programটি লক্ষ্য করো—

\begin{codebox}

\begin{english}
\ttfamily
print("First")\\
print("Second")\\
print("Third")
\end{english}

\end{codebox}

\textbf{Output:}

\begin{codebox}

\begin{english}
\ttfamily
First\\
Second\\
Third
\end{english}

\end{codebox}

Python সাধারণভাবে প্রথম Statement থেকে শুরু করে নিচের
Statementগুলো পর্যায়ক্রমে Execute করে।

\begin{center}

\textbf{Top}
\(\longrightarrow\)
\textbf{Bottom}

\end{center}

\begin{rememberbox}

Conditional Statement, Loop ও Function ব্যবহারের ক্ষেত্রে
Program-এর Execution Flow পরিবর্তিত হতে পারে।

সেগুলো পরবর্তী অংশে আলোচনা করা হবে।

\end{rememberbox}

%================================================
% Comment
%================================================

\booksubheading{Comment-এর প্রাথমিক ধারণা}

Program-এর মধ্যে Programmer-এর ব্যাখ্যা বা Note লেখার জন্য
\textbf{Comment} ব্যবহার করা হয়।

Python-এ এক লাইনের Comment সাধারণত
\textbf{\textenglish{\ttfamily \#}}
দিয়ে শুরু হয়।

\begin{examplebox}{উদাহরণ}

\begin{codebox}

\begin{english}
\ttfamily
\# This is a comment\\
print("Hello")
\end{english}

\end{codebox}

Comment Program-এর কার্যকর Instruction হিসেবে Execute হয় না।

\end{examplebox}

\begin{rememberbox}

\textbf{\color{deepOrange}MCQ তথ্য}

\textenglish{\ttfamily \#}
\(\longrightarrow\)
Single-line Comment

\end{rememberbox}

%================================================
% Program Analysis
%================================================

\booksubheading{একটি ছোট Program বিশ্লেষণ}

\begin{codebox}

\begin{english}
\ttfamily
\# My first Python program\\
print("Hello")\\
print(10 + 5)
\end{english}

\end{codebox}

\begin{tabularx}
{\linewidth}
{
|>{\RaggedRight\arraybackslash}p{0.45\linewidth}
|Y|
}

\hline

\rowcolor{softBlue}

\multicolumn{1}{|c|}{\textbf{অংশ}}
&
\multicolumn{1}{c|}{\textbf{অর্থ}}

\\

\hline

\textenglish{\ttfamily \# My first Python program}
&
Comment
\\

\hline

\textenglish{\ttfamily print("Hello")}
&
Text Output প্রদর্শন
\\

\hline

\textenglish{\ttfamily print(10 + 5)}
&
Calculation-এর Result প্রদর্শন
\\

\hline

\textenglish{\ttfamily 15}
&
দ্বিতীয়
\textenglish{\ttfamily print()}
Statement-এর Output
\\

\hline

\end{tabularx}

%================================================
% Basic Characteristics
%================================================

\booksubheading{Python Program-এর প্রাথমিক বৈশিষ্ট্য}

একটি সাধারণ Python Program-এ—

\begin{compacttwocolpoints}

\item এক বা একাধিক Statement থাকতে পারে।

\item Statement সাধারণভাবে Top to Bottom Execute হয়।

\item
\textenglish{\ttfamily print()}
ব্যবহার করে Output প্রদর্শন করা যায়।

\item Text সাধারণত Quotation Mark-এর মধ্যে লেখা হয়।

\item Calculation Expression আকারে লেখা যায়।

\item
\textenglish{\ttfamily \#}
ব্যবহার করে Single-line Comment লেখা যায়।

\item Script Program
\textenglish{\ttfamily .py}
File-এ Save করা যায়।

\end{compacttwocolpoints}

%================================================
% First Program Flow
%================================================

\booksubheading{প্রথম Program লেখার সাধারণ ধারা}

\begin{center}

\textbf{Editor}
\(\longrightarrow\)
\textbf{Python Code লেখা}
\(\longrightarrow\)
\textbf{\textenglish{\ttfamily .py} File Save}
\(\longrightarrow\)
\textbf{Run}
\(\longrightarrow\)
\textbf{Output}

\end{center}

উদাহরণ—

\begin{codebox}

\begin{english}
\ttfamily
print("My First Python Program")
\end{english}

\end{codebox}

\textbf{File:}

\begin{codebox}

\begin{english}
\ttfamily
first.py
\end{english}

\end{codebox}

\textbf{Output:}

\begin{codebox}

\begin{english}
\ttfamily
My First Python Program
\end{english}

\end{codebox}

%================================================
% Common Mistakes
%================================================

\booksubheading{সাধারণ ভুল}

\begin{compacttwocolpoints}

\item Python Code-এর সঙ্গে Output-ও Program File-এ লিখে দেওয়া।

\item
\textenglish{\ttfamily >>>}
-কে Script Program-এর অংশ মনে করা।

\item Text লেখার সময় Quotation Mark ভুলে যাওয়া।

\item
\textenglish{\ttfamily "10 + 20"}
ও
\textenglish{\ttfamily 10 + 20}
একই Result দেবে মনে করা।

\item
\textenglish{\ttfamily \#}
-কে Program Output-এর অংশ মনে করা।

\item Python Program File ভুল Extension-এ Save করা।

\item Python Program সবসময় নিচ থেকে উপরে Execute হয় মনে করা।

\end{compacttwocolpoints}

%================================================
% Final Reminder
%================================================

\begin{rememberbox}

\textbf{\color{deepOrange}মনে রাখবে}

\begin{itemize}

\item
Python Program হলো এক বা একাধিক
\textbf{Statement-এর সমষ্টি}।

\item
Programmer-এর লেখা Program Code-কে
\textbf{Source Code}
বলা হয়।

\item
Python Interpreter Source Code Execute করে।

\item
Python সাধারণভাবে Statement
\textbf{Top to Bottom}
Execute করে।

\item
\textenglish{\ttfamily print()}
ব্যবহার করে Output প্রদর্শন করা হয়।

\item
Text সাধারণত Quotation Mark-এর মধ্যে লেখা হয়।

\item
\textenglish{\ttfamily \#}
দিয়ে Single-line Comment শুরু হয়।

\item
Python Script File-এর Extension হলো
\textbf{\textenglish{\ttfamily .py}}।

\item
দ্রুত মনে রাখবে—

\begin{center}

\textbf{Source Code}
\(\longrightarrow\)
\textbf{Interpreter}
\(\longrightarrow\)
\textbf{Execution}
\(\longrightarrow\)
\textbf{Output}

\medskip

\textbf{Statement}
\(\longrightarrow\)
\textbf{Statement}
\(\longrightarrow\)
\textbf{Statement}
\(\longrightarrow\)
\textbf{Program}

\medskip

\textenglish{\ttfamily print()}
\(\longrightarrow\)
\textbf{Output}
\qquad
\textenglish{\ttfamily \#}
\(\longrightarrow\)
\textbf{Comment}

\end{center}

\end{itemize}

\end{rememberbox}

\small

%================================================
% 5.31 Python Syntax
%================================================

